\documentclass{beamer}
% \usepackage{color}
\usepackage{graphicx}
% \usepackage[all,cmtip]{xy}
\usepackage{tikz}
\usepackage{tikz-cd}
\usepackage{qtree}
\usepackage{stmaryrd}
% \usepackage{diagxy}
\usepackage{amsfonts}
\usepackage{amssymb}
\usepackage{amsmath}
\usepackage{amscd}
\usepackage{amsthm}
% \usepackage{mathpartir}
\usepackage{bbm}
%\usepackage{pictexwd,dcpic}
% \usepackage[mathcal]{euscript}
\usepackage[utf8]{inputenc}
%\usepackage[T1]{fontenc}
% \usepackage{lmodern}
\title{Explanation via Surplus Structure}
\author{Hans Halvorson}
\date{March 17, 2014}
% \usetheme{Pittsburgh}
\usetheme{Singapore}
\setbeamertemplate{navigation symbols}{}%remove navigation symbols
% \usetheme{Singapore}
\usepackage{eucal}
\begin{document}

\def\2{\mathcal}
\def\7{\mathfrak}
\begin{frame}

\maketitle

\end{frame}

\begin{frame}

\begin{itemize}
\item what is surplus structure in a scientific theory?
\item structural ockham razor (SOCK): posit only those structures that
  are needed to explain the phenomena
\item converse SOCK: if a mathematical structure is needed to explain
  the phenomena, then it represents something real
\end{itemize}

\end{frame}

\begin{frame}{quantum field theory \\ the problem of inequivalent representations}

  \textcolor{green}{OBJECTIVE:} describe a quantum system
  corresponding to a classical system with infinitely many degrees of
  freedom $q_1,p_1,q_2,p_2,\dots $

\bigskip

\textcolor{red}{OBSTACLE:} embarrassment of riches --- many
different, apparently non-equivalent, quantum descriptions




\end{frame}

\begin{frame}{inequivalent representations}

\begin{itemize}
\item quantization: given some equation $\mathbb{E}(x_1,\dots ,x_n)$,
  find a Hilbert space $H$ and operators $a_1,\dots ,a_n$ on $H$ that
  satisfy $\mathbb{E}$.
\item algebrizing: there is an abstract algebra $A$ with operators
  $a_1,\dots ,a_n\in A$ such that:
\begin{enumerate}
\item $E(a_1,\dots ,a_n)$,
\item representations $\pi :A\to B(H)$ are in one-to-one
  correspondence with quantizations of $\mathbb{E}$.
\end{enumerate}
\end{itemize}


\end{frame}

\begin{frame}{unitary equivalence}

$$
\begin{tikzcd}[ampersand replacement=\&]
A \times H \arrow{r}{\pi } \arrow{d}[left]{1_A\times w } \& H \arrow{d}[right]{w} \\
A \times K \arrow{r}[below]{\phi } \& K 
\end{tikzcd}
$$

\end{frame}


\begin{frame}{proposed solutions of INEQ}

\begin{enumerate}
\item hilbert space conservatism: there is one correct representation

  {\small $\Longrightarrow$ choosing a representation involves a
    further theoretical commitment}

\bigskip

\item algebraic imperialism: all physical information is in the
  abstract algebra of observables

  \smallskip {\small $\Longrightarrow$ choosing a representation entails no further theoretical commitment}
\end{enumerate}

\end{frame}

\begin{frame}{crash course in algebraic quantum theory}

\def\2{\mathcal}

\begin{itemize}
\item Observables are represented by operators in $A$.
\item States are defined as mappings from observables to $\mathbb{R}$.
\item A representation $\pi :A\to H$ induces a map $x\mapsto \pi
  ^*(x)$ from state vectors in $H$ to states on $\2A$.
$$ \pi ^*(x)(a) \: = \: \langle x,\pi (a)x \rangle .$$
\item GNS theorem: for every state $\omega$ of $A$, there is a Hilbert
  space $H$, a vector $x\in H$, and a representation $\pi :A\to H$
  such that 
$$\omega =\pi ^*(x) .$$
\end{itemize}


\end{frame}


\begin{frame}{problems for h.s.\ conservatism}

\begin{itemize}
\item underdetermination: what reason for choosing a representation?
\item what is gained by choosing a representation?
\item what is lost by choosing a representation?
\item what kind of epistemic commitment is it to choose a
  representation? \newline
``The system is in some state $\pi ^*(x)$.''
\end{itemize}

\end{frame}

\begin{frame}{problems for algebraic imperialism}

\begin{itemize}
\item can we do quantum theory without a Hilbert space?  \newline
% emerging theme: more structure than meets the eye 
{\small states = functions from observable algebra to numbers \newline
  dynamics = automorphism of the observable algebra \newline
  parochial observables ???}
\item too much structure? \newline
there are too many states
\end{itemize}


\end{frame}

\begin{frame}{the weak equivalence gambit}

\begin{itemize}
\item Fell's theorem: for state $x\in H$, observables $a_1,\dots
  ,a_n\in A$, and $\epsilon >0$, there is a state $y\in K$ such that 
$$ |\langle x,a_ix\rangle - \langle y,a_iy\rangle | \: < \: \epsilon .$$
\item Claim (Segal, Robinson): any two representations are empirically equivalent
\end{itemize}

\end{frame}


\begin{frame}{don't deflate representation}

\begin{itemize}

\item Ruetsche: physical equivalence $\Longrightarrow$ unitary
  equivalence

\item Fact: empirical equivalence $\Longrightarrow$ unitary
  equivalence
\end{itemize}

\end{frame}

\begin{frame}{do deflate representation}

\begin{enumerate}
\item intertranslatability $\Longrightarrow$ theoretical equivalence
  (converse of Glymour 1970)
\item Fact: there are ineq representations that are intertranslatable
\item Conclusion: there are ineq representations that are
  theoretically equivalent
\end{enumerate}


\end{frame}


% now for the free Fermi field

%% draw picture A(O) \subseteq F(O)

%% A(O) = gauge invariant elements of F(O)

%% A = inductive limit of all A(O), same with F

%% two subspaces H_{-} and H_{+} that are invariant under action of A

%% Fact: \pi _{-}(A) and \pi _{+}(A) are unitarily inequivalent, yet isomorphic

%% Puzzle: Why not just throw away one of \pi _{-} or \pi _{+} ?



\begin{frame}{free Majorana field}

\begin{itemize}
\item gauge transformation: $\gamma :F\to F$

%% puzzle: if A is invariant under \gamma , then how does \gamma change sectors?

%% perhaps Morsella explains the local endomorphisms?

\item observables are fixed points in $F$ under action of $\gamma$
\item two inequivalent representations: $\pi _+$ vacuum, $\pi _{-}$ charged
\item localized automorphisms of $A$
$$ \pi _{-}(a) \: \cong \: \pi _{+}(\theta (a)) .$$
\item $H_{+}$ and $H_{-}$ are dynamical islands
\item field operators map $H_{+}$ to $H_{-}$.
\end{itemize}


%% open question: what does the free Fermi field model explain?




\end{frame}


\begin{frame}{physical equivalence}

\begin{itemize}
\item The structure of $\pi _{+}$ is isomorphic to the structure of
  $\pi _{-}$.
\item Baker \& Halvorson: representations $\pi _{+}$ and $\pi _{-}$
  are intertranslatable.

\end{itemize}


\end{frame}

\begin{frame}{physical inequivalence}

\begin{itemize}
\item $\pi _+$ and $\pi _{-}$ are unitarily inequivalent.
\item The vacuum state $\Omega$ is in $\pi _+$ and not in $\pi
  _{-}$.
\item In fact, $\pi _{+}$ and $\pi _{-}$ are empirically inequivalent.
\end{itemize}


\end{frame}


\begin{frame}{the Buridan's ass of QFT}

\begin{enumerate}
\item The world is represented either by $\pi _{+}$ or by $\pi _{-}$, and not by both.  
%% must choose a state
%% does choosing a state require choosing a representation?
%% does choosing a state give a representation for free?
\item If the world is represented by $\pi _{+}$, it could equally well
  be represented by $\pi _{-}$.
%% the two representations are isomorphic in all ways!
%% here I agree with Baker.  These two reps have the same representational power
\end{enumerate}


\end{frame}

\begin{frame}{the dilemma restated}

\begin{enumerate}
%% Ruetsche and Earman
\item break the symmetry between the two representations
%% Baker
\item deny that there is any choice to be made
%% suppose that Buridan's ass accepts PII
%% he concludes that he is seeing double
%% .. or suffering from an illusion
%% the correct description is one in which there is just one haystack
\end{enumerate}

\bigskip
common theme: ``freedom to choose a representation'' is
theoretically insignificant

% compare with: ``freedom to choose a coordinate chart'' is
% theoretically insignificant

\end{frame}

\begin{frame}{relations between representations}

Majorana particle is its own antiparticle \newline
$\Longrightarrow$ two Majorana particle state is in the vacuum sector

\bigskip
\[ \begin{array}{lll}
\pi _{-}\otimes \pi _{-} & \cong & \pi _{+}  \\ \\
\pi _{+}\otimes \pi _{+} & \cong & \pi _{+} \end{array} \]


\end{frame}

\begin{frame}{conclusions}

\begin{enumerate}
\item the best explanation might require non-representational
  structures
\item corollary: essential $\not\Longrightarrow$ representational
\item for future research: what kinds of non-representational
  structures occur in our best theories, and what theoretical roles do
  they play?
\end{enumerate}

\end{frame}







\end{document}
